% gpr_scene.tex
% GPR acquisition geometry: antenna scanning surface, buried objects,
% hyperbolic wavefronts, soil layers.
% Colors: theme names (shibuiaccent, shibuitext, shibuiplot2, shibuiplot3)
% Bounding box: approx 10cm x 6cm
%
% Include in a frame with: % gpr_scene.tex
% GPR acquisition geometry: antenna scanning surface, buried objects,
% hyperbolic wavefronts, soil layers.
% Colors: theme names (shibuiaccent, shibuitext, shibuiplot2, shibuiplot3)
% Bounding box: approx 10cm x 6cm
%
% Include in a frame with: % gpr_scene.tex
% GPR acquisition geometry: antenna scanning surface, buried objects,
% hyperbolic wavefronts, soil layers.
% Colors: theme names (shibuiaccent, shibuitext, shibuiplot2, shibuiplot3)
% Bounding box: approx 10cm x 6cm
%
% Include in a frame with: % gpr_scene.tex
% GPR acquisition geometry: antenna scanning surface, buried objects,
% hyperbolic wavefronts, soil layers.
% Colors: theme names (shibuiaccent, shibuitext, shibuiplot2, shibuiplot3)
% Bounding box: approx 10cm x 6cm
%
% Include in a frame with: \input{tikz/gpr_scene}

\begin{tikzpicture}[
    % Base style for labels
    lbl/.style  = {font=\small, text=shibuitext},
    % Style for soil layer shading
    soil/.style = {fill=shibuicodebg, draw=shibuitext, thin},
    % Style for object
    obj/.style  = {fill=shibuiplot3!60!shibuicodebg, draw=shibuiplot3, thick},
    % Style for wave arc
    wave/.style = {draw=shibuiplot2, thin, dashed}
  ]

  % ------ coordinate definitions ------
  % Surface at y=0, depth increases downward (y negative)
  \def\surfaceY{0}
  \def\soilDepth{-3.5}    % bottom of soil shown
  \def\sceneWidth{9}      % total width in cm
  \def\antennaX{4.5}      % antenna x position

  % Object 1 (pipe, left) at (2, -1.8)
  \coordinate (obj1) at (2, -1.8);
  % Object 2 (deeper, right) at (6.5, -2.4)
  \coordinate (obj2) at (6.5, -2.4);

  % ------ soil background ------
  \fill[shibuicodebg!60!black]
    (0, \soilDepth) rectangle (\sceneWidth, \surfaceY);

  % Soil layer 1 (shallow, slightly lighter)
  \fill[shibuicodebg!80!black, opacity=0.5]
    (0, -1.0) rectangle (\sceneWidth, \surfaceY);

  % ------ surface line ------
  \draw[shibuitext, thick] (0, \surfaceY) -- (\sceneWidth, \surfaceY);
  \node[lbl, above right] at (0, \surfaceY) {surface};

  % ------ GPR antenna (moving) ------
  % Antenna body (rectangle)
  \fill[shibuiaccent!80!black]
    (\antennaX - 0.5, 0.15) rectangle (\antennaX + 0.5, 0.65);
  \draw[shibuiaccent, thick]
    (\antennaX - 0.5, 0.15) rectangle (\antennaX + 0.5, 0.65);
  % Antenna label
  \node[lbl, above] at (\antennaX, 0.65) {GPR antenna};
  % Motion arrow
  \draw[->, shibuiaccent, thick]
    (\antennaX + 0.6, 0.4) -- (\antennaX + 1.5, 0.4)
    node[lbl, right] {$x$};

  % ------ downward EM pulse arcs from antenna ------
  % Each arc is a half-ellipse opening downward, centered at the antenna.
  % Precompute radii with \pgfmathsetmacro to avoid PGF math "and" operator error.
  \foreach \r in {0.5, 1.0, 1.5} {
    \pgfmathsetmacro{\xrad}{\r / 2.5}
    \pgfmathsetmacro{\yrad}{\r / 4}
    \pgfmathsetmacro{\opac}{1 - 0.25 * \r}
    \pgfmathsetmacro{\startX}{\antennaX - \xrad}
    \draw[wave, opacity=\opac]
      (\startX, \surfaceY)
      arc[start angle=180, end angle=360,
          x radius=\xrad, y radius=\yrad];
  }

  % ------ buried objects ------
  % Object 1: circle (pipe cross-section)
  \draw[obj] (obj1) circle (0.25);
  \node[lbl, below=2pt] at (obj1) {buried pipe};

  % Object 2: smaller circle
  \draw[obj] (obj2) circle (0.18);
  \node[lbl, below=2pt] at (obj2) {cable};

  % ------ hyperbolic response on surface (B-scan schematic) ------
  % Hyperbola for object 1: vertex at x=2, t=1.4 (reflection delay)
  % Points: (x, t) where t(x) ~ sqrt((x-x0)^2/v^2 + (2d/v)^2)
  % Shown as a curve above the surface (shifted up for clarity)
  % Actually: hyperbola shown ON the B-scan panel to the right
  % Here we show the reflected wavefronts as curved arcs back to surface

  % Reflected wavefront from obj1 reaching surface
  \draw[wave, shibuiplot3, opacity=0.7]
    plot[smooth, samples=40, domain=0.3:3.7]
    (\x, {-0.05 - 0.5*(\x-2)*(\x-2)/2.5});

  % Reflected wavefront from obj2
  \draw[wave, shibuiplot3!70!white, opacity=0.5]
    plot[smooth, samples=40, domain=3.8:9.0]
    (\x, {-0.05 - 0.3*(\x-6.5)*(\x-6.5)/3.2});

  % ------ depth axis ------
  \draw[->, shibuitext, thin]
    (0, \surfaceY) -- (0, \soilDepth + 0.3)
    node[lbl, left, rotate=90, anchor=south] at (-0.25, -1.5) {depth $z$};

  % ------ B-scan panel label (right side) ------
  \node[lbl, align=center] at (8.2, -1.0)
    {$\Mimage \in \R^{\nbrows \times \nbcols}$\\[2pt]B-scan};

\end{tikzpicture}


\begin{tikzpicture}[
    % Base style for labels
    lbl/.style  = {font=\small, text=shibuitext},
    % Style for soil layer shading
    soil/.style = {fill=shibuicodebg, draw=shibuitext, thin},
    % Style for object
    obj/.style  = {fill=shibuiplot3!60!shibuicodebg, draw=shibuiplot3, thick},
    % Style for wave arc
    wave/.style = {draw=shibuiplot2, thin, dashed}
  ]

  % ------ coordinate definitions ------
  % Surface at y=0, depth increases downward (y negative)
  \def\surfaceY{0}
  \def\soilDepth{-3.5}    % bottom of soil shown
  \def\sceneWidth{9}      % total width in cm
  \def\antennaX{4.5}      % antenna x position

  % Object 1 (pipe, left) at (2, -1.8)
  \coordinate (obj1) at (2, -1.8);
  % Object 2 (deeper, right) at (6.5, -2.4)
  \coordinate (obj2) at (6.5, -2.4);

  % ------ soil background ------
  \fill[shibuicodebg!60!black]
    (0, \soilDepth) rectangle (\sceneWidth, \surfaceY);

  % Soil layer 1 (shallow, slightly lighter)
  \fill[shibuicodebg!80!black, opacity=0.5]
    (0, -1.0) rectangle (\sceneWidth, \surfaceY);

  % ------ surface line ------
  \draw[shibuitext, thick] (0, \surfaceY) -- (\sceneWidth, \surfaceY);
  \node[lbl, above right] at (0, \surfaceY) {surface};

  % ------ GPR antenna (moving) ------
  % Antenna body (rectangle)
  \fill[shibuiaccent!80!black]
    (\antennaX - 0.5, 0.15) rectangle (\antennaX + 0.5, 0.65);
  \draw[shibuiaccent, thick]
    (\antennaX - 0.5, 0.15) rectangle (\antennaX + 0.5, 0.65);
  % Antenna label
  \node[lbl, above] at (\antennaX, 0.65) {GPR antenna};
  % Motion arrow
  \draw[->, shibuiaccent, thick]
    (\antennaX + 0.6, 0.4) -- (\antennaX + 1.5, 0.4)
    node[lbl, right] {$x$};

  % ------ downward EM pulse arcs from antenna ------
  % Each arc is a half-ellipse opening downward, centered at the antenna.
  % Precompute radii with \pgfmathsetmacro to avoid PGF math "and" operator error.
  \foreach \r in {0.5, 1.0, 1.5} {
    \pgfmathsetmacro{\xrad}{\r / 2.5}
    \pgfmathsetmacro{\yrad}{\r / 4}
    \pgfmathsetmacro{\opac}{1 - 0.25 * \r}
    \pgfmathsetmacro{\startX}{\antennaX - \xrad}
    \draw[wave, opacity=\opac]
      (\startX, \surfaceY)
      arc[start angle=180, end angle=360,
          x radius=\xrad, y radius=\yrad];
  }

  % ------ buried objects ------
  % Object 1: circle (pipe cross-section)
  \draw[obj] (obj1) circle (0.25);
  \node[lbl, below=2pt] at (obj1) {buried pipe};

  % Object 2: smaller circle
  \draw[obj] (obj2) circle (0.18);
  \node[lbl, below=2pt] at (obj2) {cable};

  % ------ hyperbolic response on surface (B-scan schematic) ------
  % Hyperbola for object 1: vertex at x=2, t=1.4 (reflection delay)
  % Points: (x, t) where t(x) ~ sqrt((x-x0)^2/v^2 + (2d/v)^2)
  % Shown as a curve above the surface (shifted up for clarity)
  % Actually: hyperbola shown ON the B-scan panel to the right
  % Here we show the reflected wavefronts as curved arcs back to surface

  % Reflected wavefront from obj1 reaching surface
  \draw[wave, shibuiplot3, opacity=0.7]
    plot[smooth, samples=40, domain=0.3:3.7]
    (\x, {-0.05 - 0.5*(\x-2)*(\x-2)/2.5});

  % Reflected wavefront from obj2
  \draw[wave, shibuiplot3!70!white, opacity=0.5]
    plot[smooth, samples=40, domain=3.8:9.0]
    (\x, {-0.05 - 0.3*(\x-6.5)*(\x-6.5)/3.2});

  % ------ depth axis ------
  \draw[->, shibuitext, thin]
    (0, \surfaceY) -- (0, \soilDepth + 0.3)
    node[lbl, left, rotate=90, anchor=south] at (-0.25, -1.5) {depth $z$};

  % ------ B-scan panel label (right side) ------
  \node[lbl, align=center] at (8.2, -1.0)
    {$\Mimage \in \R^{\nbrows \times \nbcols}$\\[2pt]B-scan};

\end{tikzpicture}


\begin{tikzpicture}[
    % Base style for labels
    lbl/.style  = {font=\small, text=shibuitext},
    % Style for soil layer shading
    soil/.style = {fill=shibuicodebg, draw=shibuitext, thin},
    % Style for object
    obj/.style  = {fill=shibuiplot3!60!shibuicodebg, draw=shibuiplot3, thick},
    % Style for wave arc
    wave/.style = {draw=shibuiplot2, thin, dashed}
  ]

  % ------ coordinate definitions ------
  % Surface at y=0, depth increases downward (y negative)
  \def\surfaceY{0}
  \def\soilDepth{-3.5}    % bottom of soil shown
  \def\sceneWidth{9}      % total width in cm
  \def\antennaX{4.5}      % antenna x position

  % Object 1 (pipe, left) at (2, -1.8)
  \coordinate (obj1) at (2, -1.8);
  % Object 2 (deeper, right) at (6.5, -2.4)
  \coordinate (obj2) at (6.5, -2.4);

  % ------ soil background ------
  \fill[shibuicodebg!60!black]
    (0, \soilDepth) rectangle (\sceneWidth, \surfaceY);

  % Soil layer 1 (shallow, slightly lighter)
  \fill[shibuicodebg!80!black, opacity=0.5]
    (0, -1.0) rectangle (\sceneWidth, \surfaceY);

  % ------ surface line ------
  \draw[shibuitext, thick] (0, \surfaceY) -- (\sceneWidth, \surfaceY);
  \node[lbl, above right] at (0, \surfaceY) {surface};

  % ------ GPR antenna (moving) ------
  % Antenna body (rectangle)
  \fill[shibuiaccent!80!black]
    (\antennaX - 0.5, 0.15) rectangle (\antennaX + 0.5, 0.65);
  \draw[shibuiaccent, thick]
    (\antennaX - 0.5, 0.15) rectangle (\antennaX + 0.5, 0.65);
  % Antenna label
  \node[lbl, above] at (\antennaX, 0.65) {GPR antenna};
  % Motion arrow
  \draw[->, shibuiaccent, thick]
    (\antennaX + 0.6, 0.4) -- (\antennaX + 1.5, 0.4)
    node[lbl, right] {$x$};

  % ------ downward EM pulse arcs from antenna ------
  % Each arc is a half-ellipse opening downward, centered at the antenna.
  % Precompute radii with \pgfmathsetmacro to avoid PGF math "and" operator error.
  \foreach \r in {0.5, 1.0, 1.5} {
    \pgfmathsetmacro{\xrad}{\r / 2.5}
    \pgfmathsetmacro{\yrad}{\r / 4}
    \pgfmathsetmacro{\opac}{1 - 0.25 * \r}
    \pgfmathsetmacro{\startX}{\antennaX - \xrad}
    \draw[wave, opacity=\opac]
      (\startX, \surfaceY)
      arc[start angle=180, end angle=360,
          x radius=\xrad, y radius=\yrad];
  }

  % ------ buried objects ------
  % Object 1: circle (pipe cross-section)
  \draw[obj] (obj1) circle (0.25);
  \node[lbl, below=2pt] at (obj1) {buried pipe};

  % Object 2: smaller circle
  \draw[obj] (obj2) circle (0.18);
  \node[lbl, below=2pt] at (obj2) {cable};

  % ------ hyperbolic response on surface (B-scan schematic) ------
  % Hyperbola for object 1: vertex at x=2, t=1.4 (reflection delay)
  % Points: (x, t) where t(x) ~ sqrt((x-x0)^2/v^2 + (2d/v)^2)
  % Shown as a curve above the surface (shifted up for clarity)
  % Actually: hyperbola shown ON the B-scan panel to the right
  % Here we show the reflected wavefronts as curved arcs back to surface

  % Reflected wavefront from obj1 reaching surface
  \draw[wave, shibuiplot3, opacity=0.7]
    plot[smooth, samples=40, domain=0.3:3.7]
    (\x, {-0.05 - 0.5*(\x-2)*(\x-2)/2.5});

  % Reflected wavefront from obj2
  \draw[wave, shibuiplot3!70!white, opacity=0.5]
    plot[smooth, samples=40, domain=3.8:9.0]
    (\x, {-0.05 - 0.3*(\x-6.5)*(\x-6.5)/3.2});

  % ------ depth axis ------
  \draw[->, shibuitext, thin]
    (0, \surfaceY) -- (0, \soilDepth + 0.3)
    node[lbl, left, rotate=90, anchor=south] at (-0.25, -1.5) {depth $z$};

  % ------ B-scan panel label (right side) ------
  \node[lbl, align=center] at (8.2, -1.0)
    {$\Mimage \in \R^{\nbrows \times \nbcols}$\\[2pt]B-scan};

\end{tikzpicture}


\begin{tikzpicture}[
    % Base style for labels
    lbl/.style  = {font=\small, text=shibuitext},
    % Style for soil layer shading
    soil/.style = {fill=shibuicodebg, draw=shibuitext, thin},
    % Style for object
    obj/.style  = {fill=shibuiplot3!60!shibuicodebg, draw=shibuiplot3, thick},
    % Style for wave arc
    wave/.style = {draw=shibuiplot2, thin, dashed}
  ]

  % ------ coordinate definitions ------
  % Surface at y=0, depth increases downward (y negative)
  \def\surfaceY{0}
  \def\soilDepth{-3.5}    % bottom of soil shown
  \def\sceneWidth{9}      % total width in cm
  \def\antennaX{4.5}      % antenna x position

  % Object 1 (pipe, left) at (2, -1.8)
  \coordinate (obj1) at (2, -1.8);
  % Object 2 (deeper, right) at (6.5, -2.4)
  \coordinate (obj2) at (6.5, -2.4);

  % ------ soil background ------
  \fill[shibuicodebg!60!black]
    (0, \soilDepth) rectangle (\sceneWidth, \surfaceY);

  % Soil layer 1 (shallow, slightly lighter)
  \fill[shibuicodebg!80!black, opacity=0.5]
    (0, -1.0) rectangle (\sceneWidth, \surfaceY);

  % ------ surface line ------
  \draw[shibuitext, thick] (0, \surfaceY) -- (\sceneWidth, \surfaceY);
  \node[lbl, above right] at (0, \surfaceY) {surface};

  % ------ GPR antenna (moving) ------
  % Antenna body (rectangle)
  \fill[shibuiaccent!80!black]
    (\antennaX - 0.5, 0.15) rectangle (\antennaX + 0.5, 0.65);
  \draw[shibuiaccent, thick]
    (\antennaX - 0.5, 0.15) rectangle (\antennaX + 0.5, 0.65);
  % Antenna label
  \node[lbl, above] at (\antennaX, 0.65) {GPR antenna};
  % Motion arrow
  \draw[->, shibuiaccent, thick]
    (\antennaX + 0.6, 0.4) -- (\antennaX + 1.5, 0.4)
    node[lbl, right] {$x$};

  % ------ downward EM pulse arcs from antenna ------
  % Each arc is a half-ellipse opening downward, centered at the antenna.
  % Precompute radii with \pgfmathsetmacro to avoid PGF math "and" operator error.
  \foreach \r in {0.5, 1.0, 1.5} {
    \pgfmathsetmacro{\xrad}{\r / 2.5}
    \pgfmathsetmacro{\yrad}{\r / 4}
    \pgfmathsetmacro{\opac}{1 - 0.25 * \r}
    \pgfmathsetmacro{\startX}{\antennaX - \xrad}
    \draw[wave, opacity=\opac]
      (\startX, \surfaceY)
      arc[start angle=180, end angle=360,
          x radius=\xrad, y radius=\yrad];
  }

  % ------ buried objects ------
  % Object 1: circle (pipe cross-section)
  \draw[obj] (obj1) circle (0.25);
  \node[lbl, below=2pt] at (obj1) {buried pipe};

  % Object 2: smaller circle
  \draw[obj] (obj2) circle (0.18);
  \node[lbl, below=2pt] at (obj2) {cable};

  % ------ hyperbolic response on surface (B-scan schematic) ------
  % Hyperbola for object 1: vertex at x=2, t=1.4 (reflection delay)
  % Points: (x, t) where t(x) ~ sqrt((x-x0)^2/v^2 + (2d/v)^2)
  % Shown as a curve above the surface (shifted up for clarity)
  % Actually: hyperbola shown ON the B-scan panel to the right
  % Here we show the reflected wavefronts as curved arcs back to surface

  % Reflected wavefront from obj1 reaching surface
  \draw[wave, shibuiplot3, opacity=0.7]
    plot[smooth, samples=40, domain=0.3:3.7]
    (\x, {-0.05 - 0.5*(\x-2)*(\x-2)/2.5});

  % Reflected wavefront from obj2
  \draw[wave, shibuiplot3!70!white, opacity=0.5]
    plot[smooth, samples=40, domain=3.8:9.0]
    (\x, {-0.05 - 0.3*(\x-6.5)*(\x-6.5)/3.2});

  % ------ depth axis ------
  \draw[->, shibuitext, thin]
    (0, \surfaceY) -- (0, \soilDepth + 0.3)
    node[lbl, left, rotate=90, anchor=south] at (-0.25, -1.5) {depth $z$};

  % ------ B-scan panel label (right side) ------
  \node[lbl, align=center] at (8.2, -1.0)
    {$\Mimage \in \R^{\nbrows \times \nbcols}$\\[2pt]B-scan};

\end{tikzpicture}
